% Actual class: 01
% Slide decks: CZ3005 Module 1_Introduction; CZ3005 Module 2_Intelligent Agents and Search
% Exact scope: Module 1 pp. 1--18; Module 2 pp. 1--28, inferred from Week 2 schedule
% NTULearn supplement: None
% Human verified: Concepts discussed and checks completed; exact classroom endpoint unconfirmed

\section{第 01 节课：Intelligent Agents 与 Problem Formulation}
\label{lecture:SC3000-L01}

\lecturemeta
  {SC3000-L01}
  {2026-08-17}
  {CZ3005 Module 1；CZ3005 Module 2}
  {Module 1 pp.~1--18；Module 2 pp.~1--28（按 Week 2 schedule 推定）}
  {无}
  {知识内容已讨论并通过检查题；实际课堂终点待核对}

\subsection{本讲主线}

Artificial intelligence（人工智能）在本课程中被落实为可操作的问题：agent（智能体）如何感知 environment（环境）、依据 performance measure（性能指标）选择 rational action（理性动作），并把目标形式化为可搜索的 state space（状态空间）。历史与前沿案例只作背景，本讲的核心是 agent、task environment 与 problem formulation。

\textbf{课程提示：}Tutorial 从 Week 3 开始；两次 Lab 从 Week 6 起安排，具体信息以 NTULearn 公告为准。当前考核为 programming/lab assignment（40\%）和 final exam（60\%）。

\lecturetopic{SC3000-L01-T01}{AI 的课程视角}

\keyidea{“由机器表现出的 intelligence（智能）”只是宽泛描述；本课程用 rational agent（理性智能体）框架把 AI 转化为可定义、可计算、可评价的问题。}

AlphaGo、large language model（大语言模型）、diffusion model（扩散模型）和 scientific discovery（科学发现）展示了 AI 在特定任务上的能力，但不自动等价于 general intelligence（通用智能）或 causal understanding（因果理解）。后续课程主线是
\[
  \text{agent}\longrightarrow\text{search}\longrightarrow
  \text{decision making}\longrightarrow\text{learning}.
\]

\lecturetopic{SC3000-L01-T02}{Agent 与 Rationality}

Agent 通过 sensors（传感器）接收 percepts（感知结果），再通过 effectors / actuators（执行器）向 environment 输出 actions（动作）。Percept sequence（感知序列）是截至当前的全部感知历史。

\begin{figure}[htbp]
  \centering
  \resizebox{0.82\textwidth}{!}{\begin{tikzpicture}[
  >=Latex,
  every node/.style={font=\small},
  block/.style={draw=noteBlue,rounded corners,thick,minimum width=32mm,minimum height=14mm,align=center},
  channel/.style={draw=ntuRed,rounded corners,thick,minimum width=26mm,minimum height=10mm,align=center}
]
  \node[block] (env) {Environment\\（环境）};
  \node[channel,right=2.2cm of env,yshift=1.1cm] (sensor) {Sensors\\（传感器）};
  \node[block,right=2.2cm of sensor,yshift=-1.1cm] (agent) {Agent\\（智能体）};
  \node[channel,left=2.2cm of agent,yshift=-1.1cm] (actuator) {Actuators\\（执行器）};

  \draw[->,thick] (env) -- node[above,sloped] {percepts} (sensor);
  \draw[->,thick] (sensor) -- node[above,sloped] {$p_t$} (agent);
  \draw[->,thick] (agent) -- node[below,sloped] {$a_t$} (actuator);
  \draw[->,thick] (actuator) -- node[below,sloped] {actions} (env);
\end{tikzpicture}
}
  \caption{Agent--environment loop（智能体--环境循环）：感知流入 agent，动作反过来改变环境。}
  \label{fig:sc3000-l01-agent-loop}
\end{figure}

Rational action 是在当前 percept sequence、built-in knowledge（内置知识）和 available actions（可用动作）下，使 objective performance measure 的 expected value（期望值）最大的动作。按课件定义可形式化补充为
\[
  \boxed{a_t^*=\arg\max_{a\in\mathcal A}
  \mathbb E\!\left[M\mid p_{1:t},K,a\right]}.
\]
Rational 不等于 omniscient（全知）、必然成功或符合人类道德；判断时必须同时检查 performance measure 和 agent 当时掌握的信息。若目标函数只奖励速度，违法捷径仍可能在该错误目标下“理性”。

\textbf{对应例题：}\relatedexercise{SC3000-L01-E01}{Driverless Taxi 的任务描述与环境分类}。

\lecturetopic{SC3000-L01-T03}{Task Environment 的五组性质}

\begin{center}
\small
\begin{tabularx}{\textwidth}{@{}>{\raggedright\arraybackslash}p{3.5cm}X>{\raggedright\arraybackslash}p{4.0cm}@{}}
  \toprule
  维度 & 判断问题 & 典型例子 \\
  \midrule
  Fully / partially observable（完全/部分可观测） & Sensors 能否获得完整 state？ & Chess 完全可观测；驾驶部分可观测 \\
  Deterministic / stochastic（确定性/随机性） & 给定当前 state 与 action，下一 state 是否唯一？ & 棋规确定；真实交通含不确定性 \\
  Episodic / sequential（情节式/序列式） & 当前 action 是否影响后续决策？ & 单图分类情节式；下棋序列式 \\
  Static / dynamic（静态/动态） & Agent 思考时环境是否继续变化？ & 无计时棋局静态；驾驶动态 \\
  Discrete / continuous（离散/连续） & States、percepts 或 actions 是否连续取值？ & 棋盘离散；速度与距离连续 \\
  \bottomrule
\end{tabularx}
\end{center}

各维度彼此独立：partially observable 不等于 stochastic。Chess with a clock（带计时器的国际象棋）通常归为 fully observable、deterministic、sequential、semidynamic（半动态）和 discrete；棋盘不自行变化，但剩余时间持续减少。

\lecturetopic{SC3000-L01-T04}{Well-defined Problem Formulation}

Problem-solving agent（问题求解智能体）依次进行 goal formulation（目标制定）、problem formulation（问题形式化）、search（搜索）与 execution（执行）。一个 well-defined problem（良定义问题）可写成
\[
  \boxed{\mathcal P=(S,A,T,s_0,G,c)},
\]
其中 $S$ 是 state space，$A$ 是 action set，$T$ 是 transition model（状态转移模型），$s_0$ 是 initial state，$G$ 是 goal test，$c$ 是 step cost（单步代价）。一条路径的 path cost 为
\[
  \boxed{g(n)=\sum_{i=0}^{k-1}c(s_i,a_i,s_{i+1})}.
\]
Solution（解）是一串合法 actions，使 $s_0$ 到达满足 $G$ 的状态。评价方案时需区分 search cost（离线搜索成本）与 execution cost（在线执行成本），两者之和才是实际问题求解成本。

\textbf{严格性补充：}8-puzzle（八数码）的全部排列有 $9!$ 个，但从固定 initial state 可达的状态因 permutation parity（排列奇偶性）只有 $9!/2=181\,440$ 个。

\lecturetopic{SC3000-L01-T05}{Single-state、Belief State 与 Multiple-state Problem}

在 fully observable 且 deterministic 的模型中，agent 知道唯一 current state，这是 single-state problem（单状态问题）。若真实状态不可完全观测，agent 用 belief state（信念状态）
\[
  B_t\subseteq S
\]
表示根据当前 action--percept history 仍不能排除的 states。真实世界并非同时处于多个状态；不确定性位于 agent 的知识中。

在 deterministic 且没有新 observation 的情况下，belief update（信念更新）为
\[
  B_{t+1}=\{T(s,a_t):s\in B_t\}.
\]
Simple Vacuum World（简单吸尘器世界）有 $2\times2\times2=8$ 个 states。初始状态完全未知时，动作序列
\[
  \texttt{right},\ \texttt{suck},\ \texttt{left},\ \texttt{suck}
\]
能保证两侧都干净，其 belief states 如\cref{fig:sc3000-l01-belief-update}。

\begin{figure}[htbp]
  \centering
  \resizebox{0.97\textwidth}{!}{\begin{tikzpicture}[
  >=Latex,
  every node/.style={font=\small},
  state/.style={draw=noteBlue,rounded corners,thick,minimum width=24mm,minimum height=10mm,align=center}
]
  \node[state] (b0) {$B_0$\\$\{1,2,3,4,5,6,7,8\}$};
  \node[state,right=1.45cm of b0] (b1) {$B_1$\\$\{2,4,6,8\}$};
  \node[state,right=1.45cm of b1] (b2) {$B_2$\\$\{4,8\}$};
  \node[state,right=1.45cm of b2] (b3) {$B_3$\\$\{3,7\}$};
  \node[state,right=1.45cm of b3] (b4) {$B_4$\\$\{7\}$};

  \draw[->,thick] (b0) -- node[above] {right} (b1);
  \draw[->,thick] (b1) -- node[above] {suck} (b2);
  \draw[->,thick] (b2) -- node[above] {left} (b3);
  \draw[->,thick] (b3) -- node[above] {suck} (b4);
\end{tikzpicture}
}
  \caption{Vacuum World 的 belief-state update：动作逐步排除不可能状态，最终确定处于干净的状态 7。}
  \label{fig:sc3000-l01-belief-update}
\end{figure}

\textbf{对应例题：}\relatedexercise{SC3000-L01-E02}{Vacuum World 的保证清洁方案}。

\subsection{一分钟回顾}

Agent 在 perception--action loop（感知--行动循环）中工作；rationality 由 performance measure、已有信息与可用动作共同决定。Task environment 的五组性质必须分别判断。Problem formulation 把现实任务写成 states、actions、transition、goal test 与 cost；状态不完全可知时，以会随动作和观察更新的 belief state 表示不确定性。
